\section{Finite computation, calibration and causal composition}\label{sec:composition}
The inequalities proved so far concern Borel existence with a declared physical acquisition program. A finite source description and an executable precision certificate are additional mathematical data. We state exactly where they enter.

\begin{proposition}[Numerical implementation at a restored physical budget]\label{prop:numerics}
Suppose a chosen ideal density and dictionary at cap $\beta$ have a finite implementation whose acceptance probability, as a function of the current virtual report, differs pointwise from $h/\beta$ by at most $\eta\ge0$. Suppose its actual exported decoder values differ from the ideal exported values by at most $\delta_{\rm num}$, with no larger output width. These finite operations use the same legal acquisition program and counted live roles. Put $b=N/\kappa_+$ and
\[
 \beta=\frac{l_+b}{1+b(a_++\eta)}\ge1.
\]
If the chosen ideal pair has objective at most $\Q_m(\beta)+\xi$, the implemented program terminates almost surely at actual expected core cost at most $N$ and risk at most
\begin{equation}\label{eq:numerical-law}
 \left[
 \sqrt{\frac{u_+(1+b(a_++\eta))}{l_+}
                  (\Q_m(\beta)+\xi)+b(e_++L\eta)}
       +\delta_{\rm num}\right]^2.
\end{equation}
Here $\xi$ is a certified design/integration suboptimality, not a local optimizer's unverified gap. A uniform bound $T_{\mathrm{attempt}}$ on conditional expected physical and arithmetic time of a whole attempt yields expected time at most $bT_{\mathrm{attempt}}$. Program bits and scratch are separately those of the certified implementation.
\end{proposition}
\begin{proof}
Acceptance rounding reduces the conditional hazard by at most $\eta$. The chosen cap gives $l_+/\beta-a_+-\eta=1/b$. The same conditional-survival summation as in the main theorem yields at most $b$ attempts and cost at most $N$. Before perturbing decoder values, the one-attempt unnormalized loss changes by at most $L\eta$ under the actual joint law, because loss is bounded by $L$. Its other contribution is at most $u_+(\Q_m(\beta)+\xi)/\beta+e_+$. Summation gives the expression under the square root. At the actual committed report, the implemented and ideal exported values are within $\delta_{\rm num}$; Minkowski's inequality in the actual stopped $L^2$ law adds this distance to root-mean-square loss. It does not multiply the decoder error by the waiting time. The time claim is a separate predictable sum of the stated whole-attempt bound.
\end{proof}

In the preceding display $T_{\mathrm{attempt}}$ denotes a declared time quantity, not a free clock input to the observer. An implementation based on $m$ fixed centers in a $k$-dimensional bounded box, $p$-bit coordinate certificates and sequential squared-distance comparisons has sufficient read-only data $O(mkp)$ bits, within-call scratch $O(kp+\log m)$ bits and $O(mkp^2)$ bit operations per report, in addition to its primitive measurement and coordinate-evaluation costs. These are upper accounts, not optimal computational lower bounds. General Borel classifiers need not have such certificates. The complete dyadic implementation in the companion gives explicit rational thresholds, separate read-only/scratch counts and an independently proved profile lower.


\begin{corollary}[Finite-description matching upper under effective primitives]\label{cor:effective}
In the pure relative regular class, suppose additionally that the actual upper-program report laws are bounded above by $\bar u\nu$ on every surviving history. Assume effective compactness of the dictionary carrier and certified finite approximations of exported readouts inside $H$. Assume the calibrated report law has an effective sequence of finite physical partitions which approximates bounded finite decision vectors in $L^1(\nu)$, with certified integration of their masses and bounded task losses. Primitive classifiers and their finite costs are supplied as part of this effective interface. Then a finite-description upper with the same live and terminal role counts has risk at most a fixed multiple of $F(b,m)$ at actual expected cost at most $\kappa_+b$, for all feasible $b,m$. Every such construction has finite, separately charged program, scratch and execution accounts. No polynomial-time discovery or matching bit-complexity lower is asserted.
\end{corollary}
\begin{proof}
For $b\ge4/l_+$ choose the ideal cap $\beta=l_+b/2$. Its conditional acceptance is at least $2/b$. By Proposition~\ref{prop:regular}, its profile is at most a fixed multiple of $F(b,m)$. Represent its finite decision as functions $q_c(x)$, one for each dictionary value, whose sum is the acceptance probability. Approximate the vector on one finite physical partition by rational vectors in the subprobability simplex. The assumed effective approximation makes the sum of its $L^1(\nu)$ coordinate errors arbitrarily small, say at most $\gamma$. Domination of the actual report law bounds its actual conditional decision discrepancy by $\bar u\gamma$ on every history, and its unnormalized loss discrepancy by $L\bar u\gamma$. Approximate the actual finite exported values inside $H$ as well.

Choose $\gamma$ so small that $\bar u\gamma<1/(2b)$ and $bL\bar u\gamma\le F(b,m)$, and choose export error at most $\sqrt{F(b,m)}$. The conditional acceptance still exceeds $1/b$, so the expected cost is at most $\kappa_+b$. Summing loss and applying the stopped $L^2$ readout estimate gives a fixed multiple of $F(b,m)$. For $1\le b<4/l_+$, accept the first attempt and use one finitely described value in $H$. Bounded loss is a fixed multiple of the strictly positive lower bound for $b^{-2/\alpha}$ on this interval, uniformly in $m$.

The existence argument can be effected by finite enumeration using the supplied integration certificates and effective dictionary nets: the above choices leave strict mass and loss margins, so a candidate with certified acceptable intervals is eventually found. This need not be an efficient search. Current partition indices and rational randomization are charged within-call work; they do not create new intercall roles. All chosen descriptions and precision parameters are finite. Exact zero risk on an atom requires an exactly implementable atom classifier and readout; otherwise Proposition~\ref{prop:numerics}, not an unqualified finite-code zero claim, governs that component.
\end{proof}

\begin{proposition}[Resource-certified causal morphisms]\label{prop:composition}
A causal simulator must specify its parameter-independent strategy lift, legal command map, report transformation, internal state, clock and terminal task readout. Suppose a target program has $\ell$ live and $t$ terminal roles. A compatible simulator has at most $K$ continuing internal states, including every persistent phase, coin and clock, and at most $E$ terminal export classes sufficient for its final decoder. Then a source implementation uses at most $\ell K+tE$ state-stopped labels or $\max(\ell K,tE)$ event-stopped labels. Without a terminal sufficient-class certificate, retaining the whole simulator state gives the safe product bound.

If the complete stopped scored transcript, including mapped stopping time and premature commitment, is within total variation $\varepsilon$ of the target law uniformly over the declared models and policies, then the bounded common-task risk changes by at most $L\varepsilon$. Two compatible simulators multiply their continuing and terminal resource factors and add their deficiencies. Raw-call count and physical-time transfer require separate actual cost or killed-potential certificates. Reverse risk inequalities require a reverse simulator or an actual lower witness; terminal-law proximity alone does not provide them.
\end{proposition}
\begin{proof}
Before commitment retain the pair of target live role and continuing simulator state. After commitment retain the pair of target output role and sufficient export class. Compatibility supplies legal transitions and commands without an uncharged port or phase. Lemma~\ref{lem:roles} gives the two packing counts. Serial composition retains triples; rebracketing yields the products $K_1K_2$ and $E_1E_2$. For a loss in $[0,L]$, integration of its superlevel probabilities gives $|\E_P\ell-\E_Q\ell|\le L\|P-Q\|_{\TV}$. Total variation contracts under a common Markov kernel, so the triangle inequality adds complete-law deficiencies. Costs are generally unbounded functions of stopped paths and are not controlled by this bounded-loss argument; sum their certified increments or use Lemma~\ref{lem:drift} instead.
\end{proof}

Combining Propositions~\ref{prop:numerics} and \ref{prop:composition} adds $L\varepsilon$ to \eqref{eq:numerical-law} while charging the simulator's states and actual cost map. This composes approximation and simulation in one task metric. A bound on an allowed error is not, without attained alternatives, an unavoidable lower bound. The following construction supplies such alternatives on a single task.

\begin{theorem}[A same-task resource--calibration--deficiency law]\label{thm:joint}
Take any core class satisfying Theorem~\ref{thm:transfer} with identity exports and $B$ common live roles. Add one compulsory paid observation of an independent fair sign $E_0$, erased with probability $2\varepsilon$, $0\le\varepsilon\le1/2$, and a final paid audit of
\[
 (Z,E_0,s\delta_{\rm cal}),\qquad s\in\{-1,1\}.
\]
The fixed unknown sign $s$ affects no pre-audit law. The core, including its costs, has no access to $E_0$ beyond the initial visible tag. Use squared Euclidean loss on this product task, and charge these two additional calls. The core lower from \eqref{eq:central-lower}, evaluated with budget $N-2$ and safe projected dictionary width $M$, plus $2\varepsilon+\delta_{\rm cal}^2$, is a lower bound for either stopping interface.

Let $U(n,m)$ denote the upper expression in \eqref{eq:central-upper} for core budget $n$ and output width $m$, or its first-attempt alternative. For state stopping choose any $m\ge1$ with $1+3B+3m\le M$; for event stopping choose any $m\ge1$ with $\max(1+3B,3m)\le M$. Then an upper is
\begin{equation}\label{eq:joint-upper}
 U(N-2,m)+2\varepsilon+\delta_{\rm cal}^2.
\end{equation}
In the relative-order regular class, for $M$ large enough that these sufficient widths are a fixed fraction of $M$, this yields
\begin{equation}\label{eq:joint-law}
 R_{\mathrm{joint}}(N,M,\delta_{\rm cal},\varepsilon)
 \asymp F((N-2)/\kappa_-,M)+\delta_{\rm cal}^2+\varepsilon.
\end{equation}
The early sign channel has causal deficiency exactly $\varepsilon$ relative to the perfect-sign channel when the virtual sign must be supplied before any core call. The state counts are sufficient, not minimal for the enlarged experiment.
\end{theorem}
\begin{proof}
For the upper use one initial role, then three tagged copies of the core live roles for observed $+$, observed $-$ and erasure. Use three copies of its output dictionary. Decode the known sign or zero on erasure, and decode zero for the inaccessible last coordinate. This uses $1+3B$ live and $3m$ terminal roles, giving the stated widths by Lemma~\ref{lem:roles}. The common core upper is uniform after each tag. The other losses are exactly $2\varepsilon$ and $\delta_{\rm cal}^2$. The two compulsory calls give core budget $N-2$.

For the lower use the same actual core witness as in the central theorem. Conditioning on the initial tag does not change its full-report certificates; the initial sign and observer coins are additional independent randomization. The projected core decoder has at most $M$ values, so the actual-history proof gives its core lower. On erasure, induction over the core kernels and the stopping decisions leaves $E_0$ fair conditional on all available pre-audit information; its conditional squared prediction loss is at least one. Average the two indistinguishable fixed choices $s=\pm1$. For any last coordinate $r$, their mean loss is $\delta_{\rm cal}^2+r^2\ge\delta_{\rm cal}^2$. The core and erased-sign lower bounds hold under both of these identical pre-audit laws. Thus the three contributions hold simultaneously in one worst-case total task, rather than under incompatible separately optimized models. Regular geometry and fixed width ratios prove \eqref{eq:joint-law}.

On an erased branch every immediate virtual sign is independent of the true fair sign. The event that they agree has probability at most $1/2$, giving complete scored-law deficiency at least $\varepsilon$. Reporting the known sign otherwise and an independent fair sign on erasure gives a coupling with mismatch probability exactly $\varepsilon$. Run the same future core and audit in this coupling. The early-delivery requirement prevents waiting for the final audit to recover the sign, and proves the opposite inequality.
\end{proof}

\subsection*{A finite calibration acquisition with guarded stopping}
The reference and task-order certificates above are not inferred automatically from finite data. Here is one precise finite-sample certificate and its cost when the unknown quantities are the fixed probabilities of a physical command sampler. It also explains why high confidence alone cannot certify almost-sure stopping.

\begin{proposition}[Finite command calibration and a counted guard]\label{prop:pilot}
Suppose a physical sampler has a fixed unknown distribution $(\lambda_j)_{j=1}^J$ throughout calibration and deployment, its draws are independent, and deployment draws are independent of the current hidden input. The known target probabilities are $1/J$. Observe $n$ independent calibration draws at cost $c_{\mathrm{cal}}$ each, and let $\widehat\lambda_j$ be their frequencies. Fix $0<\zeta<1$, and choose a certified rational bound
\[
 r_n\ge\sqrt{\frac{\log(2J/\zeta)}{2n}}.
\]
Declare a relative certificate with parameter $\eta<1$ only when all intervals $[\widehat\lambda_j-r_n,\widehat\lambda_j+r_n]$ lie in $[(1-\eta)/J,(1+\eta)/J]$. The probability of a falsely declared certificate is at most $\zeta$.

On certificate failure, commit by a bounded-loss first-attempt fallback. On certificate success, deploy the corresponding upper of Theorem~\ref{thm:transfer}, with design hazard $1/b$, but force commitment with a fixed decoder after at most $T\ge1$ productive attempts. The whole-attempt cost bound $\kappa_+$ is required whether or not the certificate is true. If the true sampler lies in a tighter relative band $\eta_0<\eta$ and $2r_n\le(\eta-\eta_0)/J$, then the entire guarded procedure commits almost surely and has
\begin{align}
 \E C_{\mathrm{total}}&\le nc_{\mathrm{cal}}+\kappa_+(b+\zeta T),\label{eq:pilot-cost}\\
 R_{\mathrm{total}}&\le U+L\bigl(2\zeta+e^{-T/b}\bigr),\label{eq:pilot-risk}
\end{align}
where $U$ is the certified core upper. The calibration phase can be implemented with at most $(n+1)^J$ count-vector labels. Deployment uses a counted attempt index, multiplying its live-role bound by at most $T$; fallback and terminal roles are additional explicit roles. These resources may not be omitted from an autonomous budget.
\end{proposition}
\begin{proof}
For a Bernoulli variable $X$, the log moment-generating function of $X-\E X$ has second derivative at most $1/4$ under every exponential tilt, and value and derivative zero at the origin. Twice integrating gives $\log\E e^{t(X-\E X)}\le t^2/8$. Independence, the exponential Markov inequality and optimization in $t$ imply $\Prb(|\widehat\lambda_j-\lambda_j|>r)\le2e^{-2nr^2}$. A union bound proves that all confidence intervals contain the true probabilities except on an event of probability at most $\zeta$. Hence a false certificate has probability at most $\zeta$. If the tighter-band condition holds, this same simultaneous event puts each empirical interval within the declared wider band, so rejection of a true certificate also has probability at most $\zeta$.

On a true certificate, the conditional hazard proof bounds survival through $T$ attempts by $(1-1/b)^T\le e^{-T/b}$. The loss contributed by earlier acceptances is at most $U$; the forced last commitment adds at most $Le^{-T/b}$. Expected attempts on this branch are at most $b$. On a certificate failure the first-attempt fallback costs at most $\kappa_+$ and loses at most $L$. On a falsely accepted certificate the deterministic guard uses at most $T$ attempts and has loss at most $L$, even if the unguarded acceptance set has zero actual mass. Every attempt terminates almost surely by its finite whole-attempt cost bound. Combining these cases gives \eqref{eq:pilot-cost} and \eqref{eq:pilot-risk}, with the non-false cases bounded by $b\ge1$ attempts and the false-acceptance case by $\zeta T$ additional attempts.

A count vector stores every calibration total, including the number of draws already taken, and has at most $(n+1)^J$ possible values. At total $n$, the threshold comparison may be performed in charged scratch using a fixed certified finite arithmetic implementation; the resulting branch choice is encoded in the selected live role. The counts can then be discarded because the deployed dictionary is fixed for the nominal law and declared band, not fitted to the realized frequencies. The deployment guard requires its own retained index with $T$ values. This gives the asserted conservative finite-state accounts. Without the tighter-band assumption one still has the false-certificate probability bound and guarded termination, but no small rejection-risk bound; we do not omit that distinction.
\end{proof}

This result controls one genuinely acquired calibration quantity, not arbitrary unknown-kernel learning. It counts pilot observations, a physical fallback, the stopping guard and its memory. It prevents a rare incorrect calibration certificate from being hidden inside an unjustified robust almost-sure stopping assertion.
